\cotefile{156-5}
\scanpage{34}\mspage{33}
\diagram{\begin{tikzcd}[ampersand replacement=\&,column sep=small]
I' \arrow[d,"\varphi"'," \text{surj.}"] \& \unsure{\hookleftarrow|\leftarrow} \& J'=\varphi^{-1}(J)\arrow[d] \& \& I'^{*}\arrow[d,"{\varphi^{*}=\varphi\mid I'^{*}}"]\\
I \& \hookleftarrow \& J\ \text{partie f.} \& \text{\unsure{d'où}} \& I^{*} \& \begin{gathered}\text{\deleted{\unsure{croiss.|surj.}}}\\ \text{surjectif}\end{gathered}
\end{tikzcd}}

et on a
\diagram{\begin{tikzcd}[ampersand replacement=\&,column sep=large]
U\arrow[r,"p'","\text{surj}"']\arrow[rr,bend right=23,"p" description,"\text{surj}"'] \& I'^{*}\arrow[r,"\varphi^{*}","\text{surj}"'] \& I^{*}
\end{tikzcd}}
% slanted
\margin{\unsure{Attention}, $\varphi^{*}$ bijectif n'implique pas \unsure{néc.} $\varphi^{*}$ iso !!! Mais \unsure{si|$=$} on suppose $\varphi^{*}$ iso, ça n'implique pas a), cf plus bas.}
les \deleted{\unsure{$A$|$B$}} \unsure{strates} ouvertes de $F'$ qui sont \deleted{\unsure{dans}} \inserted{$\subset$} $U$
(i.e. pour \deleted{\unsure{dans}} $\in\unsure{G'|G_1'}$) sont les fibres de $U\xrightarrow{p'}I^{\unsure{*\alpha|**}}$,
les strates ouvertes de $F$ qui sont $\subset U$ (i.e. pour $\in G$) sont les fibres de $U\xrightarrow{p} I^{*}$.

La condition b) dit que $\varphi^{*}$ est bijective,
La condition c) dit que les \deleted{\unsure{s}}fibres de $p$ (ce sont les $A^\circ$ pour $A\in F\setminus G$) sont des fibres pour $p'$. C'est \deleted{\unsure{éq}} visiblement équivalent.

\deleted{D'autre part,}

D'autre part, que dit d) ? \unsure{Notons|Notant} \inserted{(on l'a \unsure{pr. entre|prouvé} pour $F'\ll F$)\deleted{\unsure{et là}})} que pour $A\in F$\deleted{\illegible{1}}, on a $A\in F'$ sss $A^{\circ\unsure{(F)|(F')}}$ est une strate ouverte de $F'$ (cf lemme déjà prouvé, \inserted{qu'on donnera} plus bas). Donc d) dit que pour $\forall\, A_i\in F\setminus G=I^{*}$, \deleted{\unsure{j}} la \inserted{$F$-}strate $p^{-1}(\{\deleted{\unsure{A}}i\})$\inserted{$\subset U$} est une \inserted{$F'$-}strate \deleted{\unsure{de $p$}} de $U$. Donc \unsure{cela|celà} \unsure{équiv. à visiblmt|équiv. visiblement à} $p$ bijectif.

Donc on a \[b\Longleftrightarrow c\Longleftrightarrow d\]

Montrons que c'est équiv. à a. Soit

\deleted{$F''$ comme \unsure{ds} a), d'où $I''$, \unsure{df} sur} \deleted{\unsure{$S=|F$}}

\deleted{$S=|F|=|F'|=|F''|$, on a}

\diagram{\begin{tikzcd}[ampersand replacement=\&,column sep=large,execute at end picture={\draw (\tikzcdmatrixname-1-1.north west) -- (\tikzcdmatrixname-2-1.south west); \draw (\tikzcdmatrixname-1-2.north west) -- (\tikzcdmatrixname-2-2.south east); \draw (\tikzcdmatrixname-1-3.north east) -- (\tikzcdmatrixname-2-3.south east);}]
\deleted{S}\arrow[r]\arrow[d]\arrow[rr,bend left=18,"\deleted{\varphi}"] \& \deleted{I''}\arrow[r,"\deleted{\varphi''}"]\arrow[d] \& \deleted{I}\arrow[d]\\
\deleted{|G'|=|G|=T}\arrow[r] \& \deleted{\unsure{J''|J'}}\arrow[r] \& \deleted{J}
\end{tikzcd}}
